\documentclass[11pt,a4paper]{article}
\usepackage[utf8]{inputenc}
\usepackage[italian]{babel}
\usepackage{geometry}
\usepackage{graphicx}
\usepackage{hyperref}
\usepackage{listings}
\usepackage{xcolor}
\usepackage{amsmath}
\usepackage{enumitem}
\usepackage{booktabs}
\usepackage{longtable}
\usepackage{array}
\usepackage{tcolorbox}
\usepackage{float}

\geometry{margin=2.5cm}

% Configurazione codice
\lstset{
    language=JavaScript,
    basicstyle=\ttfamily\small,
    keywordstyle=\color{blue}\bfseries,
    stringstyle=\color{red},
    commentstyle=\color{green!60!black},
    numberstyle=\tiny\color{gray},
    numbers=left,
    numbersep=5pt,
    frame=single,
    breaklines=true,
    showstringspaces=false,
    tabsize=2,
    escapeinside={(*@}{@*)}
}

% Box per evidenziare sezioni importanti
\newtcolorbox{importantbox}[1]{
    colback=red!5!white,
    colframe=red!75!black,
    title=#1,
    fonttitle=\bfseries
}

\newtcolorbox{warningbox}[1]{
    colback=orange!5!white,
    colframe=orange!75!black,
    title=#1,
    fonttitle=\bfseries
}

\newtcolorbox{tipbox}[1]{
    colback=blue!5!white,
    colframe=blue!75!black,
    title=#1,
    fonttitle=\bfseries
}

\newtcolorbox{successbox}[1]{
    colback=green!5!white,
    colframe=green!75!black,
    title=#1,
    fonttitle=\bfseries
}

\title{\textbf{WorkLogService: Stato Attuale Post-Refactoring}\\
\large Documentazione Tecnica Esaustiva - Sistema Unificato}
\author{Documentazione Sistema}
\date{\today}

\begin{document}

\maketitle
\tableofcontents
\newpage

\section{Introduzione}

\begin{importantbox}{Stato Attuale - \today}
WorkLogService è il \textbf{Cuore Operativo} del sistema di tracciamento lavoro dopo il refactoring di unificazione.
Questo componente gestisce sia log di lavoro con orari (Timer) che note di lavoro senza orari (Journal) in un'unica entità.
Se questo componente fallisce, l'app diventa inutile per l'utente.
\end{importantbox}

\subsection{Contesto Storico}

Prima del refactoring, il sistema aveva due entità separate:
\begin{itemize}
    \item \textbf{WorkLog}: Solo per tracciamento orari (startTime, endTime, durationMinutes)
    \item \textbf{WorkEntry}: Solo per note di lavoro (title, description, tags, mood)
\end{itemize}

Questo causava:
\begin{itemize}
    \item Duplicazione di codice e logica
    \item Inconsistenze nei dati (un progetto poteva avere ore ma non note)
    \item Complessità nel frontend (due API diverse)
    \item Difficoltà nel calcolare statistiche complete
\end{itemize}

Dopo il refactoring, \textbf{WorkLog} è diventato un'entità \textbf{ibrida} che supporta entrambi i casi d'uso.

\section{Architettura e Funzionamento}

\subsection{Struttura del Componente}

WorkLogService estende \texttt{BaseService}, implementando il pattern \textbf{Template Method}:

\begin{lstlisting}[caption=Costruttore WorkLogService]
class WorkLogService extends BaseService {
    constructor() {
        super(WorkLog, {
            searchFields: ['title', 'description'],
            defaultSort: { date: -1, createdAt: -1 },
            populateFields: ['projectId'],
            entityName: 'Log di lavoro',
        });
    }
}
\end{lstlisting}

\subsection{Modello Dati Unificato}

Il modello \texttt{WorkLog} supporta due modalità:

\subsubsection{Modalità Timer (con orari)}
\begin{itemize}
    \item \texttt{startTime} (HH:mm) - \textbf{Opzionale}
    \item \texttt{endTime} (HH:mm) - \textbf{Opzionale}
    \item \texttt{durationMinutes} - Calcolato automaticamente dal pre-save hook
\end{itemize}

\subsubsection{Modalità Journal (solo testo)}
\begin{itemize}
    \item \texttt{title} (required) - Titolo dell'attività
    \item \texttt{description} (optional) - Note dettagliate
    \item \texttt{tags} (optional) - Array di tag per categorizzazione
    \item \texttt{mood} (optional) - 'high', 'neutral', 'low'
    \item \texttt{durationMinutes} - 0 (default) o valore manuale
\end{itemize}

\subsubsection{Campi Comuni}
\begin{itemize}
    \item \texttt{projectId} (required) - Riferimento al progetto
    \item \texttt{date} (required, YYYY-MM-DD) - Data del lavoro
    \item \texttt{isBillable} (default: true) - Indica se fatturabile
    \item \texttt{user} - Iniettato automaticamente dal multiTenancyPlugin
\end{itemize}

\subsection{Flusso di Creazione}

\begin{figure}[H]
\centering
\begin{lstlisting}[caption=Flusso Creazione WorkLog]
[HTTP Request] POST /api/worklogs
     │
     ▼
[Controller] workLogService.create(req.tenantScope, data)
     │
     ▼
[WorkLogService.create()]
     │
     ├─► [beforeCreate()]
     │   │
     │   ├─► validateProjectOwnership()
     │   │   └─► Verifica che projectId appartenga all'utente
     │   │
     │   └─► validateTimeRange() (se orari presenti)
     │       └─► Verifica coerenza startTime/endTime
     │
     ├─► [super.create()]
     │   └─► BaseService.create() → WorkLog.create()
     │       └─► [Model pre-save hook]
     │           └─► Calcola durationMinutes se orari presenti
     │
     └─► [activityService.recordActivity()]
         ├─► WORK_SESSION_LOGGED (se orari presenti)
         └─► WORK_NOTE_CREATED (se solo testo)
\end{lstlisting}
\caption{Flusso Completo di Creazione WorkLog}
\end{figure}

\section{Interazioni con Altri Componenti}

\subsection{BaseService}

WorkLogService eredita tutte le funzionalità CRUD da BaseService:

\begin{itemize}
    \item \texttt{find(tenantScope, filters, options)} - Ricerca con filtri
    \item \texttt{findById(tenantScope, id, options)} - Ricerca per ID
    \item \texttt{create(tenantScope, data)} - Creazione (override per logica ibrida)
    \item \texttt{update(tenantScope, id, data)} - Aggiornamento
    \item \texttt{delete(tenantScope, id)} - Eliminazione
\end{itemize}

\textbf{Sicurezza}: Tutte le query includono automaticamente \texttt{\{ user: userId \}} grazie a \texttt{\_buildFilter()}.

\subsection{ActivityService}

WorkLogService registra automaticamente le attività per analytics:

\begin{lstlisting}[caption=Registrazione Attività]
// Determina tipo di attività
const activityType = (created.startTime && created.endTime) 
    ? 'WORK_SESSION_LOGGED' 
    : 'WORK_NOTE_CREATED';

await activityService.recordActivity(userId, activityType, {
    entityId: created._id,
    category: 'work',
    metadata: {
        projectId: created.projectId,
        durationMinutes,
        tags,
        timeOfDay,
        date: created.date,
        hasTimeTracking: !!(created.startTime && created.endTime),
    },
});
\end{lstlisting}

\textbf{Nota Importante}: Se la registrazione attività fallisce, il worklog viene comunque creato (non blocca l'operazione).

\subsection{Project Model}

WorkLogService verifica sempre l'ownership del progetto:

\begin{lstlisting}[caption=Validazione Ownership]
async validateProjectOwnership(tenantScope, projectId) {
    const userId = this._getUserId(tenantScope);
    
    // Query diretta con filtro esplicito
    const project = await Project.findOne({ 
        _id: projectId, 
        user: userId 
    });
    
    if (!project) {
        throw AppError.notFound('Progetto');
    }
}
\end{lstlisting}

\textbf{Sicurezza IDOR}: Previene che un utente associ log a progetti di altri utenti.

\subsection{WorkLog Model}

Il modello calcola automaticamente \texttt{durationMinutes} nel pre-save hook:

\begin{lstlisting}[caption=Pre-save Hook - Calcolo Durata]
workLogSchema.pre('save', function() {
    if (this.startTime && this.endTime) {
        const [startHour, startMin] = this.startTime.split(':').map(Number);
        const [endHour, endMin] = this.endTime.split(':').map(Number);
        
        let duration = (endHour * 60 + endMin) - (startHour * 60 + startMin);
        
        // Gestisce lavoro oltre mezzanotte
        if (duration < 0) {
            duration += 24 * 60;
        }
        
        this.durationMinutes = duration;
    } else {
        // Se orari mancano, durationMinutes = 0
        if (!this.isModified('durationMinutes')) {
            this.durationMinutes = 0;
        }
    }
});
\end{lstlisting}

\section{Funzionalità Principali}

\subsection{Metodi CRUD}

\subsubsection{create(tenantScope, data)}

Crea un nuovo worklog con logica ibrida:

\begin{itemize}
    \item \textbf{Validazione}: Ownership progetto + coerenza orari (se presenti)
    \item \textbf{Calcolo Durata}: Automatico se orari presenti, altrimenti 0
    \item \textbf{Activity Tracking}: Registra WORK\_SESSION\_LOGGED o WORK\_NOTE\_CREATED
    \item \textbf{Error Handling}: Activity tracking non blocca la creazione
\end{itemize}

\subsubsection{update(tenantScope, id, data)}

Aggiorna un worklog esistente:

\begin{itemize}
    \item \textbf{Validazione}: Se cambia projectId, verifica ownership
    \item \textbf{Validazione Orari}: Se modifica orari, richiede entrambi
    \item \textbf{Conversione}: Permette di convertire Timer → Journal (rimuovendo orari)
\end{itemize}

\subsubsection{delete(tenantScope, id)}

Elimina un worklog. \textbf{Nota}: Non elimina automaticamente le attività correlate (UserActivity rimangono per analytics storiche).

\subsection{Metodi di Query Specializzati}

\subsubsection{findByMonth(tenantScope, year, month)}

Trova tutti i log di un mese specifico:

\begin{lstlisting}[caption=Query per Mese]
async findByMonth(tenantScope, year, month) {
    const monthStr = String(month).padStart(2, '0');
    const pattern = new RegExp(`^${year}-${monthStr}-`);
    return this.find(tenantScope, { date: pattern });
}
\end{lstlisting}

\subsubsection{findByProject(tenantScope, projectId, options)}

Trova tutti i log di un progetto specifico. \textbf{Sicurezza}: Il filtro include automaticamente \texttt{user}.

\subsubsection{getWorkspaceFeed(tenantScope, filters, options)}

Feed timeline unificato con filtri avanzati:

\begin{itemize}
    \item \textbf{Filtro Progetto}: \texttt{filters.projectId}
    \item \textbf{Filtro Data}: \texttt{filters.date} (singola) o \texttt{filters.startDate/endDate} (range)
    \item \textbf{Filtro Tag}: \texttt{filters.tags} (singolo o array)
    \item \textbf{Ordinamento}: \texttt{\{ date: -1, createdAt: -1 \}} (timeline-friendly)
\end{itemize}

\textbf{Use Case}: Vista Workspace con timeline di tutte le attività (Timer + Journal).

\subsubsection{getTotalsByPeriod(tenantScope, startDate, endDate)}

Calcola totali per un periodo:

\begin{itemize}
    \item \texttt{totalMinutes} - Minuti totali lavorati
    \item \texttt{totalEarnings} - Revenue totale (ore × rate progetto)
    \item \texttt{logCount} - Numero totale di log
\end{itemize}

\textbf{Use Case}: Dashboard con statistiche mensili/annuali.

\subsubsection{getProjectStats(tenantScope, projectId)}

Statistiche dettagliate per un progetto:

\begin{itemize}
    \item \texttt{totalHours} - Ore totali lavorate
    \item \texttt{billableHours} - Ore fatturabili
    \item \texttt{nonBillableHours} - Ore non fatturabili
    \item \texttt{totalRevenue} - Revenue totale (billableHours × rate)
    \item \texttt{burnRate} - Percentuale ore consumate vs stimate (può essere null)
    \item \texttt{lastActivity} - Data ultima attività
    \item \texttt{entriesCount} - Numero totale di log/note
\end{itemize}

\textbf{Use Case}: Vista Control Room (Dashboard Progetti) con health metrics.

\section{Vantaggi per l'Utente}

\subsection{Unificazione Dati}

\begin{successbox}{Vantaggio Principale}
Un'unica entità per tracciare sia ore che note elimina la confusione e semplifica l'interfaccia utente.
\end{successbox}

\textbf{Prima del refactoring}:
\begin{itemize}
    \item L'utente doveva decidere: "È un log o una nota?"
    \item Due form diversi nel frontend
    \item Due timeline separate
    \item Statistiche incomplete (ore senza contesto o note senza ore)
\end{itemize}

\textbf{Dopo il refactoring}:
\begin{itemize}
    \item Un unico form che si adatta (mostra campi orari solo se necessario)
    \item Una timeline unificata con tutte le attività
    \item Statistiche complete (ore + note nello stesso progetto)
    \item Ricerca unificata (cerca in titoli, descrizioni, tag)
\end{itemize}

\subsection{Flessibilità d'Uso}

L'utente può:

\begin{enumerate}
    \item \textbf{Tracciare solo orari}: Inserisce startTime/endTime, il sistema calcola tutto
    \item \textbf{Scrivere solo note}: Inserisce title/description, nessun orario richiesto
    \item \textbf{Combinare entrambi}: Orari + note dettagliate per contesto completo
    \item \textbf{Convertire}: Può aggiungere/rimuovere orari in un secondo momento
\end{enumerate}

\subsection{Analytics Migliorate}

\begin{itemize}
    \item \textbf{Activity Feed}: Vede tutte le attività (Timer + Journal) in un'unica timeline
    \item \textbf{Project Health}: Metriche complete (ore + note + burn rate)
    \item \textbf{Tag System}: Categorizza attività con tag per ricerca avanzata
    \item \textbf{Mood Tracking}: Traccia umore durante il lavoro (per analytics future)
\end{itemize}

\subsection{Performance}

\begin{itemize}
    \item \textbf{Indici Ottimizzati}: Query frequenti (per mese, per progetto) sono veloci
    \item \textbf{Aggregazioni Efficaci}: Statistiche calcolate con pipeline MongoDB ottimizzate
    \item \textbf{Populate Automatico}: Progetti popolati automaticamente nelle query
\end{itemize}

\section{API Endpoints}

\subsection{GET /api/worklogs}

Lista tutti i log dell'utente (con filtri opzionali).

\textbf{Response}:
\begin{lstlisting}
{
  "success": true,
  "data": [
    {
      "id": "...",
      "projectId": {...},
      "date": "2024-01-15",
      "title": "Implementazione login",
      "startTime": "09:00",
      "endTime": "12:30",
      "durationMinutes": 210,
      "tags": ["frontend", "auth"],
      "mood": "high"
    }
  ]
}
\end{lstlisting}

\subsection{POST /api/worklogs}

Crea nuovo log (Timer o Journal).

\textbf{Body Timer}:
\begin{lstlisting}
{
  "projectId": "...",
  "date": "2024-01-15",
  "title": "Implementazione login",
  "startTime": "09:00",
  "endTime": "12:30",
  "description": "Completato sistema autenticazione",
  "tags": ["frontend", "auth"],
  "mood": "high",
  "isBillable": true
}
\end{lstlisting}

\textbf{Body Journal}:
\begin{lstlisting}
{
  "projectId": "...",
  "date": "2024-01-15",
  "title": "Note riunione",
  "description": "Discussione su nuove feature",
  "tags": ["meeting"],
  "mood": "neutral"
}
\end{lstlisting}

\subsection{GET /api/worklogs/feed}

Feed timeline con filtri avanzati.

\textbf{Query Params}:
\begin{itemize}
    \item \texttt{projectId} - Filtra per progetto
    \item \texttt{date} - Filtra per data (YYYY-MM-DD)
    \item \texttt{startDate} / \texttt{endDate} - Range date
    \item \texttt{tags} - Filtra per tag (comma-separated)
    \item \texttt{limit} - Limite risultati (default: 50)
    \item \texttt{skip} - Skip per paginazione
\end{itemize}

\subsection{GET /api/worklogs/project/:projectId/stats}

Statistiche progetto per Control Room.

\textbf{Response}:
\begin{lstlisting}
{
  "success": true,
  "data": {
    "projectId": "...",
    "projectName": "Google",
    "totalHours": 45.5,
    "billableHours": 40.0,
    "totalRevenue": 3000.00,
    "burnRate": 75.83,
    "lastActivity": "2024-01-15",
    "entriesCount": 12
  }
}
\end{lstlisting}

\section{Errori, Bug e Problemi di Logica}

\subsection{Errori Conosciuti}

\subsubsection{Activity Tracking Non Bloccante}

\begin{warningbox}{Comportamento Attuale}
Se \texttt{activityService.recordActivity()} fallisce, il worklog viene comunque creato.
L'errore viene solo loggato in console, non propagato.
\end{warningbox}

\textbf{Problema}: L'utente non viene informato se l'attività non è stata registrata per analytics.

\textbf{Impatto}: 
\begin{itemize}
    \item Analytics incomplete (manca qualche evento)
    \item Gamification non aggiornata (XP non assegnati)
    \item Streak non aggiornato
\end{itemize}

\textbf{Possibile Soluzione}: 
\begin{itemize}
    \item Retry automatico con exponential backoff
    \item Queue asincrona per activity tracking (non blocca la creazione)
    \item Notifica utente se fallisce (toast warning)
\end{itemize}

\subsubsection{Validazione Orari Parziali}

\begin{warningbox}{Edge Case}
Se l'utente invia solo \texttt{startTime} o solo \texttt{endTime}, viene lanciato un errore di validazione.
Tuttavia, se modifica un worklog esistente e imposta solo uno dei due, la validazione potrebbe non essere chiara.
\end{warningbox}

\textbf{Codice Attuale}:
\begin{lstlisting}
if (data.startTime || data.endTime) {
    if (!data.startTime || !data.endTime) {
        throw AppError.validation(
            'Se specifichi un orario, devi fornire sia startTime che endTime'
        );
    }
}
\end{lstlisting}

\textbf{Problema}: Se un worklog ha già \texttt{startTime} e l'utente aggiorna solo \texttt{endTime}, la validazione funziona. Ma se un worklog non ha orari e l'utente aggiunge solo \texttt{startTime}, fallisce.

\textbf{Impatto}: UX confusa - l'utente potrebbe non capire perché deve inserire entrambi.

\subsubsection{Calcolo Durata Oltre Mezzanotte}

\begin{tipbox}{Funzionalità Corretta}
Il sistema gestisce correttamente il lavoro oltre mezzanotte (es. 23:00 - 02:00 = 3 ore).
\end{tipbox}

\textbf{Implementazione}:
\begin{lstlisting}
let duration = (endHour * 60 + endMin) - (startHour * 60 + startMin);

if (duration < 0) {
    duration += 24 * 60; // Aggiunge 24 ore
}
\end{lstlisting}

\textbf{Nota}: Questo funziona solo se il lavoro è nella stessa giornata. Se l'utente lavora dalle 23:00 di lunedì alle 02:00 di martedì, dovrebbe creare due worklog separati (uno per giorno).

\subsection{Bug Potenziali}

\subsubsection{Race Condition in Activity Tracking}

\begin{warningbox}{Bug Potenziale}
Se due worklog vengono creati simultaneamente per lo stesso utente, l'activity tracking potrebbe avere race condition nello streak update.
\end{warningbox}

\textbf{Scenario}:
\begin{enumerate}
    \item Utente crea WorkLog A alle 23:59:59
    \item Utente crea WorkLog B alle 00:00:01 (giorno successivo)
    \item Entrambi chiamano \texttt{activityService.recordActivity()}
    \item Entrambi leggono lo stesso \texttt{user.gamification.streak}
    \item Entrambi calcolano il nuovo streak
    \item Uno sovrascrive l'altro (lost update)
\end{enumerate}

\textbf{Impatto}: Streak potrebbe non essere aggiornato correttamente.

\textbf{Possibile Soluzione}: Usare MongoDB transactions o optimistic locking.

\subsubsection{Validazione Project Ownership Race Condition}

\begin{warningbox}{Bug Potenziale}
Se un progetto viene eliminato tra la validazione e la creazione del worklog, potrebbe verificarsi un errore.
\end{warningbox}

\textbf{Scenario}:
\begin{enumerate}
    \item \texttt{beforeCreate()} valida che il progetto esiste
    \item Un altro processo elimina il progetto
    \item \texttt{super.create()} tenta di creare il worklog
    \item MongoDB lancia errore di foreign key constraint (se presente) o crea worklog orfano
\end{enumerate}

\textbf{Impatto}: Worklog con \texttt{projectId} che non esiste più.

\textbf{Possibile Soluzione}: Usare MongoDB transactions o aggiungere validazione anche in \texttt{super.create()}.

\subsection{Problemi di Logica}

\subsubsection{DurationMinutes Manuale vs Automatico}

\begin{warningbox}{Ambiguity}
Se l'utente inserisce manualmente \texttt{durationMinutes} ma fornisce anche \texttt{startTime/endTime}, quale valore ha priorità?
\end{warningbox}

\textbf{Comportamento Attuale}: Il pre-save hook calcola sempre \texttt{durationMinutes} se orari sono presenti, ignorando qualsiasi valore manuale.

\textbf{Problema}: L'utente potrebbe voler correggere manualmente la durata (es. pause non tracciate).

\textbf{Possibile Soluzione}: 
\begin{itemize}
    \item Se \texttt{durationMinutes} è esplicitamente modificato dall'utente, non sovrascriverlo
    \item Aggiungere flag \texttt{manualDuration} per indicare che è stato impostato manualmente
\end{itemize}

\subsubsection{Tag Duplicati}

\begin{warningbox}{Problema}
Non c'è validazione per tag duplicati o case-insensitive.
\end{warningbox}

\textbf{Scenario}: L'utente può inserire \texttt{["Frontend", "frontend", "FRONTEND"]} che vengono trattati come tag diversi.

\textbf{Impatto}: Ricerca per tag potrebbe non trovare tutti i risultati.

\textbf{Possibile Soluzione}: Normalizzare tag (lowercase, trim) prima del salvataggio.

\subsubsection{Billable vs Non-Billable}

\begin{tipbox}{Comportamento Attuale}
\texttt{isBillable} ha default \texttt{true}. Se l'utente crea una nota (Journal) senza orari, viene comunque marcata come billable.
\end{tipbox}

\textbf{Problema}: Una nota senza orari non dovrebbe essere billable (non ci sono ore da fatturare).

\textbf{Possibile Soluzione}: Se \texttt{durationMinutes === 0}, impostare automaticamente \texttt{isBillable = false}.

\section{Performance e Ottimizzazioni}

\subsection{Indici Database}

Il modello WorkLog ha i seguenti indici:

\begin{itemize}
    \item \texttt{\{ user: 1, date: -1 \}} - Query per mese/periodo
    \item \texttt{\{ user: 1, projectId: 1, date: -1 \}} - Query per progetto
    \item \texttt{\{ title: 'text', description: 'text' \}} - Ricerca testuale
    \item \texttt{\{ tags: 1 \}} - Ricerca per tag
\end{itemize}

\textbf{Performance}: Query frequenti (feed, stats) sono ottimizzate.

\subsection{Aggregazioni MongoDB}

Le aggregazioni per statistiche usano pipeline efficienti:

\begin{lstlisting}[caption=Esempio Aggregazione Stats]
WorkLog.aggregate([
    { $match: { user: userId, projectId: projectId } },
    {
        $group: {
            _id: null,
            totalMinutes: { $sum: '$durationMinutes' },
            billableMinutes: { $sum: { $cond: [...] } },
            entriesCount: { $sum: 1 },
            lastActivityDate: { $max: '$date' },
        },
    },
]);
\end{lstlisting}

\textbf{Performance}: Calcolo statistiche in un'unica query invece di multiple query + calcoli in memoria.

\section{Sicurezza}

\subsection{Multi-Tenancy}

\begin{successbox}{Sicurezza Garantita}
Tutte le query includono automaticamente \texttt{\{ user: userId \}} grazie a BaseService.\_buildFilter().
\end{successbox}

\textbf{Implementazione}:
\begin{lstlisting}
_buildFilter(tenantScope, additionalFilters = {}) {
    const userId = this._getUserId(tenantScope);
    return {
        user: userId,
        ...additionalFilters,
    };
}
\end{lstlisting}

\subsection{IDOR Prevention}

\begin{successbox}{Validazione Ownership}
Ogni creazione/update di worklog verifica che il progetto appartenga all'utente corrente.
\end{successbox}

\textbf{Implementazione}:
\begin{lstlisting}
async validateProjectOwnership(tenantScope, projectId) {
    const userId = this._getUserId(tenantScope);
    const project = await Project.findOne({ 
        _id: projectId, 
        user: userId 
    });
    if (!project) {
        throw AppError.notFound('Progetto');
    }
}
\end{lstlisting}

\textbf{Nota}: Il messaggio di errore è generico ("Progetto non trovato") per non rivelare se il progetto esiste ma appartiene ad altri.

\section{Testing e Manutenzione}

\subsection{Test Consigliati}

\begin{enumerate}
    \item \textbf{Unit Test}:
    \begin{itemize}
        \item Validazione ownership progetto
        \item Calcolo durata (normale e oltre mezzanotte)
        \item Validazione orari (parziali, incoerenti)
        \item Activity tracking (successo e fallimento)
    \end{itemize}
    
    \item \textbf{Integration Test}:
    \begin{itemize}
        \item Creazione worklog con progetto valido/invalido
        \item Feed con filtri multipli
        \item Statistiche progetto (con/senza estimatedHours)
        \item Race condition in activity tracking
    \end{itemize}
    
    \item \textbf{E2E Test}:
    \begin{itemize}
        \item Flusso completo: crea progetto → crea worklog → verifica stats
        \item Timeline feed con filtri
        \item Conversione Timer → Journal
    \end{itemize}
\end{enumerate}

\subsection{Manutenzione}

\begin{itemize}
    \item \textbf{Monitoraggio}: Log errori activity tracking per identificare problemi
    \item \textbf{Metriche}: Tracciare tempi di risposta query aggregate
    \item \textbf{Cleanup}: Considerare archivio worklog vecchi (> 1 anno) per performance
\end{itemize}

\section{Conclusioni}

WorkLogService è un componente \textbf{critico} e \textbf{ben progettato} che:

\begin{itemize}
    \item ✅ Unifica logica Timer e Journal in un'unica entità
    \item ✅ Garantisce sicurezza multi-tenant
    \item ✅ Fornisce API flessibili per diversi use case
    \item ✅ Ottimizza performance con indici e aggregazioni
    \item ⚠️ Ha alcuni edge case da gestire (activity tracking, validazioni)
\end{itemize}

\textbf{Raccomandazioni Future}:
\begin{enumerate}
    \item Implementare retry/queue per activity tracking
    \item Aggiungere validazione tag duplicati
    \item Gestire meglio durationMinutes manuale vs automatico
    \item Considerare transactions per operazioni critiche
    \item Aggiungere test coverage completo
\end{enumerate}

\section{Storico Modifiche}

\begin{table}[H]
\centering
\begin{tabular}{llp{8cm}}
\toprule
\textbf{Data} & \textbf{Autore} & \textbf{Modifica} \\
\midrule
2024-01-XX & Sistema & Refactoring unificazione WorkLog/WorkEntry \\
2024-01-XX & Sistema & Aggiunto supporto WORK\_NOTE\_CREATED in ActivityService \\
2024-01-XX & Sistema & Fix AppError.badRequest → AppError.validation \\
2024-01-XX & Sistema & Fix z-index modale ProjectForm \\
\today & Sistema & Documentazione completa stato attuale \\
\bottomrule
\end{tabular}
\caption{Storico Modifiche Documento}
\end{table}

\end{document}
